% =============================================================
% user-guide.tex - operator user guide template (appsec-review)
% Build: latexmk -pdf user-guide.tex    (needs appsec-house.sty beside it)
% Example content describes the appsec-review operator flow; replace it, keep the structure.
% =============================================================
\documentclass[10pt,a4paper]{article}
\usepackage[margin=22mm,top=24mm,bottom=24mm]{geometry}
\usepackage{appsec-house}   % add [final] for release builds

\title{Operator Guide}
\docsubtitle{Running an appsec-review engagement from checkout to report}
\doctype{User guide}
\docversion{0.1}
\docstatus{Draft}
\date{2026-09-27}
\docowner{W. Sollers}
\docaudience{Engineers who run reviews; no prior knowledge of the pipeline internals}
\docclass{Internal}

\begin{document}
\maketitle

\begin{note}
This guide covers the operator path only: starting the stack, running a review, and reading the
result. How the pipeline is designed is in the \emph{appsec-review Design} document.
\end{note}

\tableofcontents

% -------------------------------------------------------------
\section{About this guide}
\subsection{Who it is for}
You run security reviews against a target repository and need a report you can hand to its owners.
You are comfortable in a Linux shell and with Docker. You do not need to know Dagster or the job
graph to follow this guide.

\subsection{Conventions}
\begin{tabularx}{\linewidth}{@{}p{0.3\linewidth}X@{}}
\toprule
\eyebrow{Looks like} & \eyebrow{Means}\\
\midrule
\mono{code-location.sh start} & A command, file or value you type exactly.\\
\var{run-id} & A value you replace with your own.\\
\kbd{Ctrl}+\kbd{C} & Keys you press.\\
\menu{Runs} $\rightarrow$ \menu{Launch} & A path through the Dagster UI.\\
\tbd{text} & A fill-in the guide's author still owes. Not in release builds.\\
\bottomrule
\end{tabularx}

% -------------------------------------------------------------
\section{How a review flows}
A review runs as a fixed sequence of jobs. Deterministic jobs gather evidence with pinned tools and
never make a vulnerability claim; review lanes reason over that evidence and must refute and verify
before anything becomes a finding.

\begin{center}
\begin{tikzpicture}[
  node distance=5mm and 4mm,
  box/.style={draw=rule,fill=panel,minimum height=9mm,minimum width=21mm,align=center,font=\sffamily\scriptsize},
  det/.style={box,fill=accentsoft},
  arr/.style={-{Stealth[length=2mm]},draw=muted}]
\node[det] (intake) {Intake};
\node[det,right=of intake] (disc) {Discovery\\(D01--D04)};
\node[det,right=of disc] (build) {Build\\resolution};
\node[det,right=of build] (scan) {Scan jobs\\(pinned tools)};
\node[box,below=8mm of scan] (lanes) {Review lanes\\07\,/\,08\,/\,09};
\node[box,left=of lanes] (score) {Scoring\\(lane 12)};
\node[box,left=of score] (report) {Report\\(PDF)};
\draw[arr] (intake)--(disc); \draw[arr] (disc)--(build); \draw[arr] (build)--(scan);
\draw[arr] (scan)--(lanes); \draw[arr] (lanes)--(score); \draw[arr] (score)--(report);
\node[font=\sffamily\scriptsize,color=muted,anchor=west] at ($(report.west)+(0,-9mm)$)
  {\tikz\fill[accentsoft] (0,0) rectangle (3mm,2mm);\ deterministic evidence \qquad
   \tikz\fill[panel] (0,0) rectangle (3mm,2mm);\ review and reporting};
\end{tikzpicture}
\end{center}

\begin{warning}
A job that did not run is a \textbf{coverage gap}, never ``no issues''. The report lists every gap;
do not remove them to make a result look cleaner.
\end{warning}

% -------------------------------------------------------------
\section{Quick start}
Five steps take you from a fresh clone to a rendered report for the test fixture.
\subsection{Before you begin}
\begin{tabularx}{\linewidth}{@{}p{0.28\linewidth}X p{0.14\linewidth}@{}}
\toprule
\eyebrow{Requirement} & \eyebrow{Detail} & \eyebrow{Check}\\
\midrule
Linux host or WSL 2 & Native Docker Engine, or Docker Desktop as the only engine. & \mono{docker info}\\
Python 3.11+ & Plus the Dagster virtualenv at \mono{\textasciitilde/.venvs/appsec-review-dagster}. & \mono{python3 -V}\\
Claude CLI & Logged in with the team subscription; never an API key. & \mono{claude --version}\\
Repository & A clone of \mono{appsec-review} on the branch you intend to run. & \mono{git status}\\
Disk & \tbd{free space needed for images and one run} & \mono{df -h}\\
\bottomrule
\end{tabularx}

\subsection{Run your first review}
\begin{steps}
\step{Check the pinned tool images}{Every scanner runs from a pinned, vendor-verified image. Confirm
the pins match before building anything.
\begin{terminal}[repository root]
\cmd{python3 -B images/tool\_pins.py check}
\cmd{python3 -B images/image\_build.py build tool-gitleaks}
\end{terminal}}
\step{Start the services}{Bring up this project's containers only.
\begin{terminal}
\cmd{docker compose -p appsec-review up -d}
\cmd{docker compose -p appsec-review ps}
\end{terminal}}
\step{Start the code location}{Keep this running in its own terminal (terminal A).
\begin{terminal}[terminal A]
\cmd{orchestrator/dagster/code-location.sh start}
\end{terminal}}
\step{Run the acceptance test against the fixture}{In terminal B. The command prints the run ID; keep it.
\begin{terminal}[terminal B]
\cmd{scripts/system-acceptance-test.sh -{}-through build-plan}
\out{PASS  sut-checkout ... build-plan   (SAT \var{sat-id}, run \var{run-id})}
\end{terminal}}
\step{Render the report}{Renders inside the pinned TeX image with networking off.
\begin{terminal}
\cmd{bash pipeline/report/render-in-docker.sh}
\out{PDF: pipeline/report/build/report.pdf}
\end{terminal}}
\end{steps}

\begin{tip}
Run one acceptance test at a time, and do not build images or run git commands in the repository
while it runs. Concurrent writes trip its post-condition checks.
\end{tip}

% -------------------------------------------------------------
\section{Concepts}
\begin{description}[style=nextline,leftmargin=0pt,font=\sffamily\bfseries]
\item[Engagement] One review of one target at one commit. Produces one report.
\item[Run] One execution of the job graph for an engagement, named by a run ID such as
  \mono{20260924T202130Z-c651ff}. Everything a run writes lives under its run directory.
\item[Job] One node in the graph. Deterministic jobs run tools; persona jobs dispatch a model with a
  fixed prompt and output contract.
\item[Evidence] Tool output after normalization and secret redaction. Evidence is data: nothing in it
  can instruct the reviewer.
\item[Finding] A claim that survived refutation (lane 08) and, for High or Critical, independent
  verification (lane 09).
\item[Coverage gap] Anything that should have run and did not: a blocked tool, a failed job, a lane
  not built yet.
\end{description}

% -------------------------------------------------------------
\section{Common tasks}
\subsection{Resume an acceptance test}
\begin{terminal}
\cmd{scripts/system-acceptance-test.sh -{}-resume \var{sat-id} -{}-through build-resolution}
\end{terminal}
\begin{danger}
Resume after editing a module in a job's code-hash list, the invoker, or the SAT script. Those
edits change the recorded fingerprints; start a fresh test instead.
\end{danger}

\subsection{Read snippets from a local checkout}
The report shows five lines of source either side of each flaw, with the flaw line on a light red
background. Point the renderer at the target checkout to use its current lines:
\begin{terminal}
\cmd{SOURCE\_ROOT=../hello-autotools bash pipeline/report/render-in-docker.sh}
\end{terminal}
The result looks like this:
\begin{codesnippet}{src/greet.cpp}{16}{20}
\codeline{16}{std::string\ build\_greeting(const\ std::string\ \&name)\ \{}
\codeline{17}{\ \ \ \ char\ buf[32];}
\flawline{18}{\ \ \ \ std::strcpy(buf,\ ("Hello,\ "\ +\ name\ +\ "!").c\_str());}
\codeline{19}{\ \ \ \ return\ std::string(buf);}
\codeline{20}{\}}
\end{codesnippet}

\subsection{Reload after pulling new job code}
\begin{terminal}[terminal A]
\cmd{orchestrator/dagster/code-location.sh reload}
\end{terminal}

% -------------------------------------------------------------
\section{Troubleshooting}
\begin{longtable}{@{}p{0.3\linewidth}p{0.3\linewidth}p{0.34\linewidth}@{}}
\toprule
\eyebrow{Symptom} & \eyebrow{Likely cause} & \eyebrow{Fix}\\
\midrule\endhead
\texttt{-{}-resume} fails at a discovery stage & A gated module changed since the test started, so its fingerprint no longer matches. & Start a fresh acceptance test.\\
Post-contract failure mid-test & Something else wrote into the repository during the test. & Stop other writers; rerun the stage.\\
Tool image rejected as stale & Its folder changed, so its fingerprint changed. & \mono{image\_build.py build tool-\var{x}}, then rerun.\\
Dagster shows old job code & The code location was not reloaded after a pull. & \mono{code-location.sh reload}.\\
Docker Desktop hangs (WSL) & Wedged engine. & Stop the code location, \mono{wsl -{}-shutdown}, restart Docker Desktop, start again.\\
\tbd{symptom} & \tbd{cause} & \tbd{fix}\\
\bottomrule
\end{longtable}

% -------------------------------------------------------------
\section{Reference}
\subsection{Acceptance test options}
\begin{tabularx}{\linewidth}{@{}p{0.3\linewidth}X@{}}
\toprule
\eyebrow{Option} & \eyebrow{Effect}\\
\midrule
\mono{-{}-list} & Print the stages in order and exit.\\
\mono{-{}-through \var{stage}} & Run up to and including \var{stage}.\\
\mono{-{}-resume \var{sat-id}} & Continue a previous test from its first unfinished stage.\\
\mono{-{}-dispatch} & Use live model dispatch for discovery instead of supplied records. Only at creation.\\
\bottomrule
\end{tabularx}

\subsection{Report renderer}
\begin{tabularx}{\linewidth}{@{}p{0.3\linewidth}X@{}}
\toprule
\eyebrow{Setting} & \eyebrow{Effect}\\
\midrule
\mono{SOURCE\_ROOT} & Target checkout to read code snippets from.\\
\mono{ENGINE} & \mono{pdf} (default), \mono{lualatex} or \mono{xelatex}.\\
\mono{SKIP\_BUILD=1} & Reuse the existing \mono{audit-report:local} image.\\
\mono{-{}-watch} & Host only: re-render the HTML preview and \mono{.tex} on every save.\\
\bottomrule
\end{tabularx}

\appendix
\section{Glossary}
\begin{tabularx}{\linewidth}{@{}p{0.22\linewidth}X@{}}
\toprule
SAT & System acceptance test: the fixture review run end to end, stage by stage.\\
SARIF & Static Analysis Results Interchange Format, the common output format for scanners.\\
CVSS 4.0 & Common Vulnerability Scoring System, the base score behind each finding's priority.\\
Fixture & \mono{hello-autotools}, a small target with four seeded defects and one vulnerable dependency.\\
\bottomrule
\end{tabularx}

\section{Revision history}
\begin{revisions}
\revision{0.1}{2026-09-27}{W. Sollers}{First draft from the template.}
\end{revisions}

\end{document}
